\clearpage
\begin{examplebox}

\textbf{\textcolor{CExample}{先选最快方法（高考常见题型，以下均为自编例题）}}
\begin{itemize}[leftmargin=*,itemsep=0.35em,topsep=0.3em]
  \item 同底比较大小：\textbf{先看底数与 $1$ 的关系，再比较指数}。
  \item 底数不同：先尝试化成同底，不能化同底时寻找中间值 $1$ 等。
  \item 图像与值域：\textbf{先找过哪一点、整条图像在水平线哪边}。
\end{itemize}

\medskip
\textbf{\textcolor{CExample}{例 1（快记转快用：同底、底数大于 $1$）}}
\begin{center}\begin{tikzpicture}[font=\small,>=stealth,line join=round]
\draw[gray!40,rounded corners=3pt] (0,0) rectangle (14.2,8.3);\node at (7.1,7.700000000000001) {例1：沿同一条递增曲线，从左到右高度增加};
\begin{scope}[shift={(1.8058823529411765,2.11)},x=5.88235294117647cm,y=2.3cm]\draw[->,gray] (-0.12,0)--(0.9,0) node[right]{$x$};\draw[->,gray] (0,-0.2)--(0,1.8) node[above]{$y$};\draw[blue!75!black,thick,solid,domain=-0.05:0.8,samples=180] plot(\x,{pow(1.7,\x)});\fill[orange!85!black] (0.3,1.172558924272542) circle(2.2pt);\fill[orange!85!black] (0.5,1.3038404810405297) circle(2.2pt);\node[below=5pt] at (0.3,0){$0.3$};\node[below=5pt] at (0.5,0){$0.5$};\draw[gray,densely dotted] (0.3,0)--(0.3,1.172558924272542);\draw[gray,densely dotted] (0.5,0)--(0.5,1.3038404810405297);\end{scope}\node[anchor=west,align=left,text width=5.0cm] at (8.4,4.0) {$a=1.7>1$\\[9pt]$0.3<0.5$\\[9pt]$1.7^{0.3}<1.7^{0.5}$\\[9pt]比较指数，不用求小数幂};
\node[text width=12.6cm,align=center] at (7.1,.65) {固定的是底数1.7；横坐标表示指数，左点低、右点高。};
\end{tikzpicture}\end{center}

\textbf{观察：}先看底数1.7大于1，再看0.3在0.5左边。\quad\textbf{原因：}同一条递增曲线上，左点低于右点。

\textbf{考试使用：}免求小数，直接写比较结果。



\textbf{题目：}比较 $1.7^{0.3}$ 与 $1.7^{0.5}$ 的大小。

\textbf{最快解法：}$1.7>1$，函数随指数增大而增大；因为 $0.3<0.5$，所以
\[
  \boxed{1.7^{0.3}<1.7^{0.5}}.
\]
\textbf{为什么用它最快：}不用分别求两次小数幂。

\medskip

\end{examplebox}
\clearpage
\begin{examplebox}
\textbf{\textcolor{CExample}{例 2（同底、底数小于 $1$）}}
\begin{center}\begin{tikzpicture}[font=\small,>=stealth,line join=round]
\draw[gray!40,rounded corners=3pt] (0,0) rectangle (14.2,8.3);\node at (7.1,7.700000000000001) {例2：底数小于1，右边的点反而更低};
\begin{scope}[shift={(1.82,2.161111111111111)},x=2.4cm,y=2.5555555555555554cm]\draw[->,gray] (-0.3,0)--(2.2,0) node[right]{$x$};\draw[->,gray] (0,-0.2)--(0,1.6) node[above]{$y$};\draw[orange!80!black,thick,dashed,domain=-0.2:2,samples=180] plot(\x,{pow(.5,\x)});\fill[orange!85!black] (0.8,0.5743491774985174) circle(2.2pt);\fill[orange!85!black] (1.2,0.43527528164806206) circle(2.2pt);\node[below=5pt] at (0.8,0){$0.8$};\node[below=5pt] at (1.2,0){$1.2$};\draw[gray,densely dotted] (0.8,0)--(0.8,0.5743491774985174);\draw[gray,densely dotted] (1.2,0)--(1.2,0.43527528164806206);\end{scope}\node[anchor=west,align=left,text width=5.0cm] at (8.4,4.0) {图取$a=\frac12$\\[9pt]$0.8<1.2$\\[9pt]$a^{0.8}>a^{1.2}$\\[9pt]对任意$0<a<1$同样成立};
\node[text width=12.6cm,align=center] at (7.1,.65) {比较指数前先判断底数；下降的曲线把左右顺序变成相反的高低顺序。};
\end{tikzpicture}\end{center}

\textbf{观察：}左点高、右点低。\quad\textbf{原因：}乘0与1之间的正数会变小。

\textbf{考试使用：}先判底数，不等号与指数比较相反。



\textbf{题目：}已知 $0<a<1$，比较 $a^{0.8}$ 与 $a^{1.2}$。

\textbf{最快解法：}底数小于 $1$ 时，指数越大，值越小；因为 $0.8<1.2$，所以
\[
  \boxed{a^{0.8}>a^{1.2}}.
\]
\textbf{理解检查：}想成 $a=\frac12$，每多乘一次 $\frac12$ 就会变小。

\medskip

\end{examplebox}
\clearpage
\begin{examplebox}
\textbf{\textcolor{CExample}{例 3（解指数不等式：先化同底）}}
\begin{center}\begin{tikzpicture}[font=\small,>=stealth,line join=round]
\draw[gray!40,rounded corners=3pt] (0,0) rectangle (14.2,8.3);\node at (7.1,7.700000000000001) {例3：交点横坐标2，把大于关系分成左右两侧};
\begin{scope}[shift={(3.1999999999999997,2.1177966101694916)},x=1.5cm,y=0.7796610169491526cm]\draw[->,gray] (-1.4,0)--(2.6,0) node[right]{$x$};\draw[->,gray] (0,-0.6)--(0,5.3) node[above]{$y$};\draw[blue!75!black,thick,solid,domain=-1.1:2.2,samples=180] plot(\x,{pow(2,\x)});\draw[orange!80!black,thick,dashed,domain=-1.1:2.12,samples=180] plot(\x,{pow(4,\x-1)});\fill[orange!85!black] (2,4) circle(2.2pt);\node[below left=5pt] at (0,0){$0$};\node[below=5pt] at (2,0){$2$};\draw[gray,densely dotted] (2,0)--(2,4);\end{scope}\node[anchor=west,align=left,text width=5.0cm] at (8.4,4.0) {蓝：$2^x$\\[8pt]橙：$4^{x-1}=2^{2x-2}$\\[8pt]交点$(2,4)$\\[8pt]要蓝线高于橙线\\[8pt]得到$x<2$，不含交点};
\node[text width=12.6cm,align=center] at (7.1,.65) {化成同底后比较x与2x-2；图像帮助核对解集在交点左边。};
\end{tikzpicture}\end{center}

\textbf{观察：}交点左边蓝线高，右边橙线高。\quad\textbf{原因：}化同底后比较x与2x-2。

\textbf{考试使用：}严格大于不取交点；用x=0回验。



\textbf{题目：}解不等式 $2^{x}>4^{x-1}$。

\textbf{最快解法：}先把 $4$ 化成 $2^2$，再由 $2>1$ 保持不等号方向：
\[
  2^x>2^{2x-2}\quad\Longleftrightarrow\quad x>2x-2
  \quad\Longleftrightarrow\quad\boxed{x<2}.
\]
\textbf{检查：}取 $x=0$，左边 $1$、右边 $1/4$，确实符合。

\medskip

\end{examplebox}
\clearpage
\begin{examplebox}
\textbf{\textcolor{CExample}{例 4（底数有字母：一定分类）}}
\begin{center}\begin{tikzpicture}[font=\small,>=stealth,line join=round]
\draw[gray!40,rounded corners=3pt] (0,0) rectangle (14.2,8.3);\node at (7.1,7.700000000000001) {例4：先看指数在横轴哪边，再看底数决定方向};
\begin{scope}[shift={(2.3,3.9499999999999997)},x=1.2cm,y=0.7666666666666666cm]\draw[->,gray] (-1,0)--(4,0) node[right]{$x$};\draw[->,gray] (0,-3)--(0,3) node[above]{$u$};\draw[blue!75!black,thick,solid,domain=-0.65:3.65,samples=180] plot(\x,{\x*(\x-3)});\fill[orange!85!black] (0,0) circle(2.2pt);\fill[orange!85!black] (3,0) circle(2.2pt);\fill[orange!85!black] (1.5,-2.25) circle(2.2pt);\node[below left=5pt] at (0,0){$0$};\node[below=5pt] at (3,0){$3$};\node[above=5pt] at (1.5,-2.25){$u<0$};\node[above=6pt] at (3.55,1.9525){$u>0$};\end{scope}\node[anchor=west,align=left,text width=5.0cm] at (8.4,4.0) {令$u=x^2-3x$\\[8pt]$u=0$在$x=0,3$\\[8pt]$u>0$：$x<0$或$x>3$\\[8pt]$u<0$：$0<x<3$\\[8pt]$a>1$选$u>0$\\[5pt]$0<a<1$选$u<0$};
\node[text width=12.6cm,align=center] at (7.1,.65) {要a的u次方大于1：底数大于1时选正指数，底数小于1时选负指数。};
\end{tikzpicture}\end{center}

\textbf{观察：}抛物线在横轴上方还是下方。\quad\textbf{原因：}与1比较等于和a的0次方比较。

\textbf{考试使用：}先按底数分类，再选指数正负区间。



\textbf{题目：}已知 $a>0,\ a\neq1$，解 $a^{x^2-3x}>1$。

\textbf{解题关键：}把 $1$ 写成 $a^0$，分两种情况：
\[
\begin{aligned}
 a>1:&\quad x^2-3x>0
 \quad\Longleftrightarrow\quad\boxed{x<0\ \text{或}\ x>3};\\
 0<a<1:&\quad x^2-3x<0
 \quad\Longleftrightarrow\quad\boxed{0<x<3}.
\end{aligned}
\]
\textbf{为什么要分类：}同样的指数，底数在 $1$ 的两边，不等号方向相反。$x=0,3$ 时左边恰等于 $1$，故均不取。

\medskip

\end{examplebox}
\clearpage
\begin{examplebox}
\textbf{\textcolor{CExample}{例 5（判断图像与对称点）}}
\begin{center}\begin{tikzpicture}[font=\small,>=stealth,line join=round]
\draw[gray!40,rounded corners=3pt] (0,0) rectangle (14.2,8.3);\node at (7.1,7.700000000000001) {关于y轴对称：横坐标变号，高度和距离都不变};
\begin{scope}[shift={(4.1,2.161111111111111)},x=1.0714285714285714cm,y=0.7301587301587301cm]\draw[->,gray] (-2.8,0)--(2.8,0) node[right]{$x$};\draw[->,gray] (0,-0.7)--(0,5.6) node[above]{$y$};\draw[blue!75!black,thick,solid,domain=-2.3:2.15,samples=180] plot(\x,{pow(2,\x)});\draw[orange!80!black,thick,dashed,domain=-2.3:2.15,samples=180] plot(\x,{pow(2,-\x)});\fill[orange!85!black] (-2,4) circle(2.2pt);\fill[orange!85!black] (2,4) circle(2.2pt);\fill[orange!85!black] (0,1) circle(2.2pt);\node[above=7pt] at (-1.25,4.65){$(-2,4)$};\node[above=7pt] at (1.25,4.65){$(2,4)$};\draw[<->,gray] (-2,4)--(0,4);\draw[<->,gray] (0,4)--(2,4);\node[below=4pt] at (-1,4){$2$};\node[below=4pt] at (1,4){$2$};\node[below=5pt] at (-2,0){$-2$};\node[below left=5pt] at (0,0){$0$};\node[below=5pt] at (2,0){$2$};\end{scope}\node[anchor=west,align=left,text width=5.0cm] at (8.4,4.0) {蓝：$f(x)=2^x$\\[8pt]橙：$g(x)=2^{-x}=f(-x)$\\[8pt]$(2,4)\leftrightarrow(-2,4)$\\[8pt]一般为$(t,f(t))$\\[4pt]$\leftrightarrow(-t,f(t))$};
\node[text width=12.6cm,align=center] at (7.1,.65) {同一水平线上的两点距y轴相等；从一条曲线可直接翻出另一条。};
\end{tikzpicture}\end{center}

\textbf{观察：}同一水平线的两点左右等距。\quad\textbf{原因：}倒数底数等于指数变号。

\textbf{考试使用：}对应点只改横坐标符号，交点为(0,1)。



\textbf{题目：}已知 $(2,4)$ 在 $y=2^x$ 的图像上。求它在 $y=(\frac12)^x$ 上对应的对称点，并求这两条图像的交点。

\textbf{最快解法：}两图关于 $y$ 轴对称，所以 $(2,4)$ 对应\textbf{$(-2,4)$}；图像相交时
\[
  2^x=2^{-x}\quad\Longleftrightarrow\quad x=-x
  \quad\Longleftrightarrow\quad x=0,
\]
交点为\textbf{$(0,1)$}。不必画精确曲线。

\medskip

\end{examplebox}
\clearpage
\begin{examplebox}
\textbf{\textcolor{CExample}{例 6（变式：图像上下左右平移）}}
\begin{center}\begin{tikzpicture}[font=\small,>=stealth,line join=round]
\draw[gray!40,rounded corners=3pt] (0,0) rectangle (14.2,8.3);\node at (7.1,7.700000000000001) {例6：共同点和平行靠近的水平线一起平移};
\begin{scope}[shift={(4.2875,2.079333333333333)},x=0.9375cm,y=0.6133333333333333cm]\draw[->,gray] (-3.4,0)--(3,0) node[right]{$x$};\draw[->,gray] (0,-0.7)--(0,6.8) node[above]{$y$};\draw[blue!75!black,thick,solid,domain=-3:1.6,samples=180] plot(\x,{pow(2,\x)});\draw[violet,thick,dashed,domain=-3:2.6,samples=180] plot(\x,{pow(2,\x-1)+3});\fill[orange!85!black] (0,1) circle(2.2pt);\fill[orange!85!black] (1,4) circle(2.2pt);\draw[violet,densely dotted] (-3.1,3)--(2.8,3);\node[above=6pt] at (-2.8,3){$3$};\draw[->,gray] (0,1)--(1,1)--(1,4);\node[below left=5pt] at (0,0){$0$};\node[below=5pt] at (1,0){$1$};\end{scope}\node[anchor=west,align=left,text width=5.0cm] at (8.4,4.0) {蓝：$y=2^x$\\[8pt]紫：$y=2^{x-1}+3$\\[8pt]$(0,1)\to(1,4)$\\[8pt]右移1，再上移3\\[8pt]$y=0\to y=3$\\[8pt]新值域$(3,+\infty)$};
\node[text width=12.6cm,align=center] at (7.1,.65) {每个点的横坐标加1、纵坐标加3；原来取不到的0也随之变成取不到的3。};
\end{tikzpicture}\end{center}

\textbf{观察：}点与水平线怎样一起移动。\quad\textbf{原因：}横坐标加1、纵坐标加3。

\textbf{考试使用：}先令指数为0找点，再由恒正确定值域。



\textbf{题目：}求 $y=2^{x-1}+3$ 的值域、图像经过的一个容易确定的点，以及它会无限接近的水平直线。

\textbf{最快解法：}因为 $2^{x-1}>0$，所以 $y>3$；令 $x=1$，得 $y=4$。因此
\[
  \boxed{\text{值域 }(3,+\infty),\quad\text{过点 }(1,4),\quad\text{水平渐近线 }y=3}.
\]
\textbf{关联旧知识：}原图 $y=2^x$ 向右移 $1$ 格、向上移 $3$ 格；它原来靠近的 $y=0$ 也向上移成 $y=3$。

\medskip
\end{examplebox}
\clearpage
\begin{examplebox}
\textbf{\textcolor{CExample}{举一反三（自检）}}
\begin{center}\begin{tikzpicture}[font=\small,>=stealth,line join=round]
\draw[gray!40,rounded corners=3pt] (0,0) rectangle (14.2,8.3);\node at (7.1,7.700000000000001) {自检1：底数小于1，大于等于转成指数小于等于};
\begin{scope}[shift={(2.492857142857143,3.1136363636363633)},x=1.0714285714285714cm,y=0.41818181818181815cm]\draw[->,gray] (-1.3,0)--(4.3,0) node[right]{$x$};\draw[->,gray] (0,-3.5)--(0,7.5) node[above]{$u$};\draw[teal,thick,solid,domain=-1:4,samples=180] plot(\x,{2*\x-1});\draw[violet,thick,dashed,domain=-1:4,samples=180] plot(\x,{\x+2});\fill[orange!85!black] (3,5) circle(2.2pt);\node[above left=5pt] at (0,0){$0$};\node[below=5pt] at (3,0){$3$};\draw[gray,densely dotted] (3,0)--(3,5);\end{scope}\node[anchor=west,align=left,text width=5.0cm] at (8.4,4.0) {画两个指数：\\[5pt]青$2x-1$，紫$x+2$\\[8pt]原底数为$\frac13<1$\\[8pt]要青线不高于紫线\\[8pt]$2x-1\leq x+2$\\[8pt]解集$x\leq3$};
\node[text width=12.6cm,align=center] at (7.1,.65) {这里只画指数的两条直线；先利用指数函数递减翻转比较方向，再从交点读解集。};
\end{tikzpicture}\end{center}

\textbf{观察：}两条指数直线在哪一侧满足大小关系。\quad\textbf{原因：}原函数底数小于1，比较方向翻转。

\textbf{考试使用：}小于等于保留交点，得到x小于等于3。



（1）$\left(\frac13\right)^{2x-1}\geq\left(\frac13\right)^{x+2}$ 的解集是什么？
\textbf{提示：}小于 $1$ 的底数会把不等号翻转。\textbf{答案：}$x\leq3$。

\end{examplebox}
\clearpage
\begin{examplebox}

\begin{center}\begin{tikzpicture}[font=\small,>=stealth,line join=round]
\draw[gray!40,rounded corners=3pt] (0,0) rectangle (14.2,8.3);\node at (7.1,7.700000000000001) {自检2：指数为0的位置给出点，整体减1给出水平线};
\begin{scope}[shift={(5.42,2.8328571428571427)},x=1.2cm,y=0.6571428571428571cm]\draw[->,gray] (-3.6,0)--(1.4,0) node[right]{$x$};\draw[->,gray] (0,-1.8)--(0,5.2) node[above]{$y$};\draw[orange!80!black,thick,dashed,domain=-3.2:1,samples=180] plot(\x,{pow(.25,\x+2)-1});\fill[orange!85!black] (-2,0) circle(2.2pt);\draw[gray,densely dotted] (-3.3,-1)--(1.1,-1);\node[below=6pt] at (-1,-1){$y=-1$};\node[below=5pt] at (-2,0){$-2$};\node[above left=5pt] at (0,0){$0$};\end{scope}\node[anchor=west,align=left,text width=5.0cm] at (8.4,4.0) {$y=(\frac14)^{x+2}-1$\\[8pt]$x=-2$：指数为0\\[8pt]过$(-2,0)$\\[8pt]整体向下移1\\[8pt]水平渐近线$y=-1$\\[8pt]值域$(-1,+\infty)$};
\node[text width=12.6cm,align=center] at (7.1,.65) {底数小于1，靠近水平线的方向向右；y=-1仍取不到。};
\end{tikzpicture}\end{center}

\textbf{观察：}经过(-2,0)，向右靠近y=-1。\quad\textbf{原因：}正的指数项整体减1后仍大于-1。

\textbf{考试使用：}不要把新渐近线继续写成y=0。

（2）函数 $y=(\frac14)^{x+2}-1$ 经过哪个容易确定的点？值域是什么？
\textbf{提示：}先令指数为 $0$。\textbf{答案：}过 $(-2,0)$，值域 $(-1,+\infty)$。

\end{examplebox}
